\documentclass[11pt,a4paper]{report}

% ─── Encodage & langue ───────────────────────────────────────
\usepackage[utf8]{inputenc}
\usepackage[T1]{fontenc}
\usepackage[french]{babel}
\usepackage{lmodern}

% ─── Mise en page ────────────────────────────────────────────
\usepackage[margin=2.3cm]{geometry}
\usepackage{microtype}
\usepackage{setspace}
\onehalfspacing

% ─── Couleurs ────────────────────────────────────────────────
\usepackage{xcolor}
\definecolor{navy}{HTML}{0F172A}
\definecolor{teal}{HTML}{14B8A6}
\definecolor{tealdark}{HTML}{0E7C72}
\definecolor{lightgray}{HTML}{F1F5F9}
\definecolor{midgray}{HTML}{64748B}
\definecolor{taken}{HTML}{16A34A}
\definecolor{late}{HTML}{CA8A04}
\definecolor{missed}{HTML}{DC2626}
\definecolor{codebg}{HTML}{F8FAFC}

% ─── Titres ──────────────────────────────────────────────────
\usepackage{titlesec}
\titleformat{\chapter}[display]
  {\bfseries\color{navy}}
  {\Large\color{teal}\MakeUppercase{\chaptertitlename} \thechapter}
  {0.4em}
  {\Huge}
\titlespacing*{\chapter}{0pt}{0pt}{1.4ex}
\titleformat{\section}{\Large\bfseries\color{navy}}{\thesection}{0.6em}{}
\titleformat{\subsection}{\large\bfseries\color{tealdark}}{\thesubsection}{0.6em}{}
\titlespacing*{\section}{0pt}{2.2ex plus 1ex minus .2ex}{1ex}

% ─── Tableaux ────────────────────────────────────────────────
\usepackage{tabularx}
\usepackage{booktabs}
\usepackage{array}
\usepackage{colortbl}
\renewcommand{\arraystretch}{1.35}
\newcolumntype{L}[1]{>{\raggedright\arraybackslash}p{#1}}

% ─── Listes ──────────────────────────────────────────────────
\usepackage{enumitem}
\setlist{leftmargin=1.4em,topsep=2pt,itemsep=1pt}
\usepackage{pifont}

% ─── Boîtes / encadrés ───────────────────────────────────────
\usepackage[most]{tcolorbox}
\newtcolorbox{infobox}[1][]{colback=lightgray,colframe=teal,boxrule=1pt,
  arc=2pt,left=8pt,right=8pt,top=6pt,bottom=6pt,#1}
\newtcolorbox{warnbox}[1][]{colback=missed!6,colframe=missed,boxrule=1pt,
  arc=2pt,left=8pt,right=8pt,top=6pt,bottom=6pt,#1}
\newtcolorbox{stepbox}[1][]{colback=teal!6,colframe=tealdark,boxrule=1pt,
  arc=2pt,left=8pt,right=8pt,top=6pt,bottom=6pt,#1}

% ─── Code ────────────────────────────────────────────────────
\usepackage{listings}
\lstset{
  basicstyle=\ttfamily\footnotesize,
  backgroundcolor=\color{codebg},
  frame=single,rulecolor=\color{midgray!40},
  breaklines=true,columns=fullflexible,
  xleftmargin=4pt,xrightmargin=4pt,
  aboveskip=8pt,belowskip=8pt,
}

% ─── Organigramme (TikZ) ─────────────────────────────────────
\usepackage{tikz}
\usetikzlibrary{shapes.geometric, arrows.meta, positioning, calc}
\tikzset{
  startstop/.style={rectangle, rounded corners=10pt, draw=tealdark, very thick,
    fill=teal!18, text width=8.4cm, align=center, minimum height=1cm,
    font=\small\bfseries},
  process/.style={rectangle, draw=navy, thick, fill=lightgray,
    text width=8.4cm, align=center, minimum height=0.95cm, font=\small},
  decision/.style={diamond, aspect=2.4, draw=late, very thick, fill=late!12,
    text width=4.6cm, align=center, inner sep=1pt, font=\small\bfseries},
  io/.style={trapezium, trapezium left angle=70, trapezium right angle=110,
    draw=tealdark, thick, fill=teal!8, text width=7.6cm, align=center,
    minimum height=0.9cm, font=\small},
  flowline/.style={-{Stealth[length=2.6mm]}, thick, draw=midgray},
}

% ─── Liens ───────────────────────────────────────────────────
\usepackage{hyperref}
\hypersetup{colorlinks=true,linkcolor=tealdark,urlcolor=tealdark,
  pdftitle={DisrtuCare - Chapitre 3},pdfauthor={Projet DisrtuCare}}

\usepackage{fancyhdr}
\pagestyle{fancy}\fancyhf{}
\fancyhead[L]{\small\color{midgray}DisrtuCare}
\fancyhead[R]{\small\color{midgray}Chapitre 3 --- Fonctionnement \& d\'eveloppement logiciel}
\fancyfoot[C]{\small\color{midgray}\thepage}
\renewcommand{\headrulewidth}{0.4pt}

% ═════════════════════════════════════════════════════════════
\begin{document}

\setcounter{chapter}{2}
\chapter{Fonctionnement du syst\`eme et d\'eveloppement logiciel}

\vspace{-1ex}
{\color{teal}\rule{\linewidth}{1.5pt}}
\vspace{1ex}

Ce chapitre d\'ecrit comment le syst\`eme \textbf{DisrtuCare} fonctionne r\'eellement,
une fois assembl\'e. Apr\`es avoir pr\'esent\'e les composants mat\'eriels au chapitre
pr\'ec\'edent, nous expliquons ici \textbf{ce qui se passe quand le syst\`eme est en
marche}~: comment le distributeur d\'elivre une pilule \`a l'heure pr\'evue, comment
l'application mobile a \'et\'e con\c cue, comment les deux communiquent, et enfin comment
toute la cha\^ine se d\'eroule \og du r\'eglage de l'heure par l'utilisateur jusqu'\`a la
pilule dans le plateau \fg. Le propos reste volontairement \textbf{accessible}~: les
notions informatiques sont expliqu\'ees par analogie plut\^ot que par le d\'etail du code.

\begin{infobox}
\textbf{Vue en une phrase.} DisrtuCare associe un \textbf{distributeur autonome} (un petit
ordinateur \og NodeMCU \fg\ qui pilote un moteur, une horloge et une LED) \`a une
\textbf{application mobile compagnon} qui sert de tableau de bord, de carnet de suivi et de
t\'el\'ecommande. Les deux dialoguent par \textbf{WiFi}, sur le r\'eseau domestique du patient.
\end{infobox}

% ═════════════════════════════════════════════════════════════
\section{Fonctionnement du syst\`eme mat\'eriel}

\subsection{Le principe g\'en\'eral}

Le c\oe ur du distributeur est un \textbf{NodeMCU V3}, un microcontr\^oleur de la famille
ESP8266. C'est l'\'equivalent d'un tr\`es petit ordinateur~: il ex\'ecute en boucle un
programme (le \emph{firmware}) et int\`egre nativement le \textbf{WiFi}, ce qui lui permet
de communiquer avec le t\'el\'ephone sans c\^able ni module suppl\'ementaire. Autour de lui
s'articulent quatre organes~:

\begin{tabularx}{\linewidth}{L{3.6cm} X}
\toprule
\rowcolor{lightgray}\textbf{Organe} & \textbf{R\^ole physiologique du syst\`eme} \\
\midrule
Moteur pas \`a pas \textbf{28BYJ-48} (via le pilote \textbf{ULN2003}) & Le \og muscle \fg~: il fait tourner le carrousel de pilules d'un compartiment, lib\'erant la dose. \\
Horloge temps r\'eel \textbf{DS3231} & La \og montre \fg~: elle garde l'heure exacte, m\^eme apr\`es une coupure de courant, gr\^ace \`a une pile bouton. \\
\textbf{LED} bleue & Le \og signal d'alerte \fg~: elle clignote pour signaler qu'une dose vient d'\^etre d\'elivr\'ee et attend confirmation. \\
\textbf{Bouton} poussoir & Le \og retour patient \fg~: en appuyant, l'utilisateur confirme qu'il a bien pris sa pilule. \\
\'Ecran \textbf{LCD} 16$\times$2 & L'\og afficheur \fg~: il montre l'heure, l'horaire et l'adresse de connexion au d\'emarrage. \\
\bottomrule
\end{tabularx}

\subsection{Le moteur~: comment une pilule est lib\'er\'ee}

La distribution repose enti\`erement sur la \textbf{rotation pr\'ecise} du moteur. Le
28BYJ-48 est un \emph{moteur pas \`a pas}~: au lieu de tourner librement, il avance par
petits \og pas \fg\ commandables un par un, ce qui permet de ma\^itriser exactement de
combien il tourne. Dans DisrtuCare~:

\begin{itemize}
  \item Un tour complet du moteur correspond \`a \textbf{4096 demi-pas}.
  \item Le carrousel poss\`ede \textbf{4 compartiments}~; un compartiment correspond donc
        \`a $4096 \div 4 = \mathbf{1024}$ \textbf{demi-pas}.
  \item Chaque demi-pas dure 2 millisecondes, soit \textbf{environ 2 secondes} de rotation
        pour amener le compartiment suivant au-dessus de l'ouverture.
  \item D\`es la rotation termin\'ee, le programme \textbf{coupe l'alimentation des
        bobines} du moteur, pour \'eviter qu'il chauffe inutilement \`a l'arr\^et.
\end{itemize}

\begin{infobox}
\textbf{Pourquoi un moteur \og pas \`a pas \fg ?} Parce qu'on a besoin d'une pr\'ecision
m\'edicale~: il faut d\'elivrer \emph{exactement une} dose, ni plus ni moins. Compter les
pas garantit que le carrousel s'arr\^ete pile sur le compartiment suivant.
\end{infobox}

\subsection{L'horloge~: savoir quand distribuer}

Pour distribuer \og \`a 8h00 \fg, le syst\`eme doit conna\^itre l'heure. C'est le r\^ole
du module \textbf{DS3231}, une horloge aliment\'ee par une petite pile qui continue de
compter le temps m\^eme appareil \'eteint. Le firmware a \'et\'e con\c cu avec une
\textbf{s\'ecurit\'e intelligente}~: si l'horloge mat\'erielle n'est pas pr\'esente, il
bascule sur une \textbf{horloge logicielle de secours}, mise \`a l'heure par le t\'el\'ephone
lors de la connexion. Le syst\`eme reste donc fonctionnel m\^eme sans le module DS3231.

\begin{infobox}
Tant que l'heure n'est pas connue de fa\c con fiable, le programme \textbf{interdit toute
distribution} (fonction \texttt{rtcIsRunning}). Cela \'evite qu'une dose parte par erreur
\`a minuit, heure par d\'efaut d'un appareil non r\'egl\'e.
\end{infobox}

\subsection{La boucle de fonctionnement permanente}

Une fois allum\'e, le distributeur ex\'ecute sans arr\^et une \textbf{boucle} qui
encha\^ine quatre v\'erifications, plusieurs fois par seconde et \textbf{sans jamais se
bloquer}~:

\begin{enumerate}
  \item \textbf{R\'epondre au t\'el\'ephone}~: traiter les \'eventuelles demandes de l'application.
  \item \textbf{V\'erifier l'horaire}~: l'heure actuelle correspond-elle \`a la dose du matin ou du soir~?
  \item \textbf{Surveiller le bouton}~: l'utilisateur vient-il de confirmer une prise~?
  \item \textbf{Animer la LED}~: la faire clignoter tant qu'une confirmation est attendue.
\end{enumerate}

\noindent C'est ce fonctionnement \textbf{autonome} qui fait la robustesse du dispositif~:
m\^eme si le t\'el\'ephone est \'eteint ou hors r\'eseau, le distributeur conna\^it l'heure,
d\'elivre les doses programm\'ees et attend la confirmation du patient.

\subsection{La m\'emoire permanente (EEPROM)}

L'horaire (heure du matin, heure du soir) et le nom du m\'edicament sont \'ecrits dans une
\textbf{m\'emoire non volatile} appel\'ee EEPROM. Concr\`etement, ces r\'eglages
\textbf{survivent \`a une coupure de courant}~: apr\`es un red\'emarrage, le distributeur
retrouve son horaire sans avoir besoin du t\'el\'ephone. Un \og octet t\'emoin \fg\ permet
au programme de savoir si la m\'emoire a d\'ej\`a \'et\'e initialis\'ee.

% ═════════════════════════════════════════════════════════════
\section{D\'eveloppement logiciel}

Le logiciel de DisrtuCare se compose de \textbf{deux programmes distincts} qui coop\`erent~:
le \textbf{firmware embarqu\'e} (d\'ecrit ci-dessus, \'ecrit en langage Arduino~/~C++) et
l'\textbf{application mobile}. Cette section d\'etaille la conception de l'application.

\subsection{La technologie choisie et pourquoi}

L'application a \'et\'e d\'evelopp\'ee avec \textbf{Expo / React Native}, une technologie
qui permet d'\'ecrire \textbf{une seule fois} le code et de le faire fonctionner aussi bien
sur \textbf{Android} que sur \textbf{iOS}. Le langage utilis\'e est le \textbf{TypeScript}.
Ce choix \'evite de d\'evelopper deux applications s\'epar\'ees et acc\'el\`ere donc
nettement la r\'ealisation.

\begin{tabularx}{\linewidth}{L{4.2cm} X}
\toprule
\rowcolor{lightgray}\textbf{Brique logicielle} & \textbf{\`A quoi elle sert} \\
\midrule
\textbf{Expo / React Native} & Construire l'interface et la faire tourner sur t\'el\'ephone. \\
\textbf{expo-sqlite} & Une base de donn\'ees locale, dans le t\'el\'ephone~: l'app fonctionne \textbf{hors-ligne}, sans serveur ni internet. \\
\textbf{expo-notifications} & Programmer des \textbf{rappels} (notifications) aux heures de prise. \\
\textbf{react-native-reanimated} & Rendre les animations fluides et agr\'eables. \\
\bottomrule
\end{tabularx}

\subsection{L'organisation de l'application en \'ecrans}

L'interface est organis\'ee en \textbf{quatre onglets}, pens\'es pour une utilisation simple,
notamment par des \textbf{personnes \^ag\'ees}~:

\begin{tabularx}{\linewidth}{L{2.8cm} X}
\toprule
\rowcolor{lightgray}\textbf{\'Ecran} & \textbf{Fonction} \\
\midrule
\textbf{Tableau de bord} & Affiche la \textbf{prochaine dose} (avec compte \`a rebours), un \textbf{anneau d'observance} sur 30 jours, et les doses du jour (matin~/~soir) avec leur statut. \\
\textbf{Historique} & Journal jour par jour des prises, avec un code couleur. \\
\textbf{Appareil} & Permet de \textbf{se connecter au distributeur} (saisie de son adresse), et offre des boutons de test (moteur, LED, synchronisation). \\
\textbf{R\'eglages} & D\'efinit les \textbf{heures} matin et soir, le \textbf{nom du m\'edicament}, et active~/~d\'esactive les notifications. \\
\bottomrule
\end{tabularx}

\subsection{La base de donn\'ees locale}

Toutes les informations sont conserv\'ees \textbf{dans le t\'el\'ephone}, dans une petite
base de donn\'ees SQLite. Elle comporte deux tables~:

\begin{itemize}
  \item \textbf{\texttt{schedule}} (les r\'eglages)~: une seule ligne contenant le nom du
        m\'edicament, les heures du matin et du soir, l'\'etat des notifications et
        l'adresse du distributeur.
  \item \textbf{\texttt{dose\_logs}} (le carnet de prises)~: une entr\'ee par dose, avec la
        date, le type (matin~/~soir) et le statut~:
        {\color{taken}\textbf{prise}} / {\color{late}\textbf{en retard}} /
        {\color{missed}\textbf{manqu\'ee}}. Une r\`egle garantit \textbf{une seule entr\'ee
        matin et une seule entr\'ee soir par jour}.
\end{itemize}

\subsection{Les rappels (notifications)}

\`A chaque enregistrement de l'horaire, l'application \textbf{reprogramme automatiquement
quatre rappels quotidiens}~: un rappel le matin, un le soir, et deux \textbf{relances 30
minutes plus tard} si la prise n'a pas \'et\'e confirm\'ee. Ces rappels sont \textbf{locaux}~:
ils s'affichent sur le t\'el\'ephone m\^eme sans connexion internet.

\subsection{L'accessibilit\'e~: une exigence de conception}

Le public cible \'etant les \textbf{patients \^ag\'es}, l'interface a \'et\'e con\c cue avec~:
grandes polices (18 \`a 22 px), \textbf{fort contraste} (fond bleu marine, accent turquoise),
\textbf{grandes zones tactiles} (au moins 56 px) et un \textbf{code couleur intuitif}
(vert~/~jaune~/~rouge) pour le statut de chaque dose.

% ═════════════════════════════════════════════════════════════
\section{Organigramme (logigramme) du fonctionnement}

L'organigramme ci-dessous r\'esume le \textbf{cycle complet d'une dose}, depuis la
v\'erification de l'heure par le distributeur jusqu'\`a l'enregistrement de la prise dans
l'application. Les losanges repr\'esentent des \textbf{d\'ecisions}~; les rectangles, des
\textbf{actions}.

\vspace{1ex}
\begin{center}
\begin{tikzpicture}[node distance=8mm and 0mm]

\node (start)  [startstop] {D\'emarrage du distributeur \\ (lecture de l'heure et de l'horaire en m\'emoire)};
\node (clock)  [decision, below=of start] {L'heure $=$ heure d'une dose~?};
\node (wait)   [process, right=14mm of clock, text width=3.2cm] {Attendre \\ (boucle)};
\node (motor)  [process, below=14mm of clock] {Le moteur tourne d'un compartiment \\ $\rightarrow$ la pilule tombe dans le plateau};
\node (led)    [process, below=of motor] {La LED clignote~; un \'ev\'enement \\ \og dose d\'elivr\'ee \fg\ est enregistr\'e};
\node (app)    [io, below=of led] {L'application (qui interroge le distributeur \\ toutes les 3 s) affiche l'\'ev\'enement};
\node (btn)    [decision, below=of app] {Le patient appuie \\ sur le bouton~?};
\node (notif)  [process, right=10mm of btn, text width=3.6cm] {Rappel renvoy\'e \\ apr\`es 30 min};
\node (conf)   [process, below=16mm of btn] {LED \'eteinte~; \'ev\'enement \\ \og prise confirm\'ee \fg};
\node (log)    [process, below=of conf] {L'app enregistre la dose comme \\ \textbf{\og prise \fg} et met \`a jour l'observance};
\node (end)    [startstop, below=of log] {Dose termin\'ee \\ (remise \`a z\'ero le lendemain \`a minuit)};

\draw [flowline] (start) -- (clock);
\draw [flowline] (clock) -- node[above, font=\footnotesize\bfseries, text=missed] {non} (wait);
\draw [flowline] (wait.north) |- ([yshift=3mm]start.south east) -| (clock.north east);
\draw [flowline] (clock) -- node[right, font=\footnotesize\bfseries, text=taken] {oui} (motor);
\draw [flowline] (motor) -- (led);
\draw [flowline] (led)   -- (app);
\draw [flowline] (app)   -- (btn);
\draw [flowline] (btn) -- node[above, font=\footnotesize\bfseries, text=missed] {non} (notif);
\draw [flowline] (notif.south) |- (app.east);
\draw [flowline] (btn) -- node[right, font=\footnotesize\bfseries, text=taken] {oui} (conf);
\draw [flowline] (conf) -- (log);
\draw [flowline] (log)  -- (end);

\end{tikzpicture}
\end{center}

\begin{infobox}
\textbf{Lecture de l'organigramme.} Tant que l'heure ne correspond pas, le distributeur
\og tourne en rond \fg\ (boucle d'attente). Au moment voulu, il agit~: il distribue,
signale, et attend la confirmation du patient. Si le patient n'appuie pas, l'application le
relance~; quand il appuie, la prise est enregistr\'ee automatiquement.
\end{infobox}

% ═════════════════════════════════════════════════════════════
\section{Int\'egration mat\'eriel--logiciel}

C'est le point cl\'e du projet~: faire \textbf{dialoguer} le distributeur et le t\'el\'ephone.

\subsection{Le moyen de communication~: le WiFi}

Le distributeur et le t\'el\'ephone sont connect\'es au \textbf{m\^eme r\'eseau WiFi
domestique}. Le NodeMCU se comporte alors comme un \textbf{petit serveur web}~: il poss\`ede
une \textbf{adresse} (par exemple \texttt{192.168.11.109}) que l'utilisateur saisit dans
l'onglet \textbf{Appareil} de l'application. Cette adresse s'affiche sur l'\'ecran LCD au
d\'emarrage, puis elle est \textbf{m\'emoris\'ee} par l'application pour les fois suivantes.

\subsection{Le langage commun~: des messages HTTP/JSON}

Les deux appareils \'echangent de courts messages au format \textbf{JSON} (un format texte
simple et lisible) via le protocole \textbf{HTTP}, le m\^eme que celui des sites web. Cinq
\og adresses \fg\ (endpoints) sont d\'efinies sur le distributeur~:

\begin{tabularx}{\linewidth}{L{2.5cm} X}
\toprule
\rowcolor{lightgray}\textbf{Message} & \textbf{Signification} \\
\midrule
\texttt{/status}  & \og Donne-moi ton \'etat \fg\ (heure, horaire, dernier \'ev\'enement). \\
\texttt{/sync}    & \og Voici le nouvel horaire \fg\ (matin, soir, m\'edicament). \\
\texttt{/settime} & \og Voici l'heure exacte \fg\ (pour r\'egler l'horloge). \\
\texttt{/led}     & \og Allume ou \'eteins la LED \fg. \\
\texttt{/motor}   & \og Fais tourner le moteur \fg\ (bouton de test). \\
\bottomrule
\end{tabularx}

\subsection{Le sens distributeur $\rightarrow$ application~: le \og sondage \fg}

Un microcontr\^oleur ne peut pas facilement \og appeler \fg\ le t\'el\'ephone de
lui-m\^eme. La solution retenue est simple et robuste~: c'est l'application qui
\textbf{interroge le distributeur toutes les 3 secondes} (m\'ecanisme dit de \emph{polling}).
Chaque r\'eponse contient un \textbf{num\'ero d'\'ev\'enement}. Quand ce num\'ero change,
l'application comprend qu'un \'ev\'enement vient de se produire (une dose d\'elivr\'ee, une
prise confirm\'ee) et r\'eagit en cons\'equence.

\begin{infobox}
\textbf{Une analogie.} Plut\^ot que le distributeur \og t\'el\'ephone \fg\ \`a l'app,
c'est l'app qui demande r\'eguli\`erement~: \og Du nouveau~? \fg. Tant que le num\'ero
d'\'ev\'enement est le m\^eme, rien n'a chang\'e~; d\`es qu'il augmente, l'app sait qu'un
fait nouveau s'est produit.
\end{infobox}

\subsection{Le pipeline complet~: du r\'eglage \`a la prise}

Voici le sc\'enario de bout en bout, qui relie tout ce qui pr\'ec\`ede~:

\begin{stepbox}
\textbf{1. R\'eglage (utilisateur).} Dans l'\'ecran R\'eglages, l'utilisateur d\'efinit
08:00 / 20:00. L'application~: \emph{(a)} enregistre l'horaire dans sa base locale~;
\emph{(b)} l'envoie au distributeur via \texttt{/sync} (qui le sauve en EEPROM et l'affiche
sur le LCD)~; \emph{(c)} programme les notifications de rappel.
\end{stepbox}

\begin{stepbox}
\textbf{2. Distribution (mat\'eriel).} \`A 08:00, le distributeur d\'etecte que l'heure
correspond \`a la dose du matin. Le moteur tourne d'un compartiment, la pilule tombe dans le
plateau, la LED se met \`a clignoter, et un \'ev\'enement \og dose d\'elivr\'ee \fg\ est
enregistr\'e.
\end{stepbox}

\begin{stepbox}
\textbf{3. Notification (logiciel).} Au m\^eme moment, le t\'el\'ephone affiche son rappel
local \og C'est l'heure de votre dose du matin \fg. En interrogeant le distributeur,
l'application voit le nouvel \'ev\'enement et l'affiche dans son journal.
\end{stepbox}

\begin{stepbox}
\textbf{4. Confirmation (utilisateur + mat\'eriel).} Le patient prend sa pilule et
\textbf{appuie sur le bouton} du distributeur. La LED s'\'eteint, et un \'ev\'enement \og
prise confirm\'ee \fg\ est enregistr\'e.
\end{stepbox}

\begin{stepbox}
\textbf{5. Tra\c cabilit\'e (logiciel).} L'application capte la confirmation et
\textbf{enregistre automatiquement la dose comme \og prise \fg} dans le carnet, sans aucune
action manuelle. L'\textbf{anneau d'observance} (30 jours) se recalcule.
\end{stepbox}

\begin{infobox}
\textbf{Le point fort de l'int\'egration}~: un simple \textbf{geste physique} sur le
distributeur (appuyer sur le bouton) \textbf{met \`a jour le dossier num\'erique} du patient.
Le patient n'a rien \`a saisir~; le suivi d'observance se construit tout seul.
\end{infobox}

% ═════════════════════════════════════════════════════════════
\section{Difficult\'es rencontr\'ees et solutions apport\'ees}

La r\'ealisation a soulev\'e plusieurs difficult\'es concr\`etes, mat\'erielles comme
logicielles. Les principales sont r\'esum\'ees ci-dessous.

\begin{tabularx}{\linewidth}{L{5.1cm} X}
\toprule
\rowcolor{lightgray}\textbf{Difficult\'e} & \textbf{Solution apport\'ee} \\
\midrule
\textbf{Le moteur ne tournait pas} malgr\'e un c\^ablage en apparence correct. &
Les fils de commande \'etaient branch\'es sur les broches D1--D4 alors que le programme
pilotait D5--D8. Le \textbf{d\'eplacement des fils sur D5--D8} a r\'esolu le probl\`eme~;
ces broches \'evitent en plus les broches \og sensibles au d\'emarrage \fg\ de l'ESP8266. \\

\textbf{Alimentation insuffisante du moteur.} Le moteur \'etait reli\'e \`a une broche
n'ayant pas la tension n\'ecessaire. & Reraccordement sur la broche \textbf{5V (VU)} reli\'ee
directement \`a l'USB, et installation du \textbf{cavalier d'alimentation} du pilote ULN2003. \\

\textbf{Coupures du microcontr\^oleur} (\og brownout \fg) au d\'emarrage du moteur, qui
tirait trop de courant via le port USB du PC. & D\'ebrancher l'alimentation moteur pendant
le t\'el\'eversement~; la solution d\'efinitive est une \textbf{alimentation 5V s\'epar\'ee}
pour le moteur, avec une masse commune. \\

\textbf{Impossible de conna\^itre l'adresse du distributeur} sans \'ecran. &
Le firmware \textbf{affiche l'adresse sur le LCD} pendant 5 secondes au d\'emarrage (et
l'imprime aussi sur le port s\'erie pour le d\'ebogage). \\

\textbf{Risque de distribution parasite} si l'heure n'est pas connue. & Le programme
\textbf{interdit toute distribution tant que l'horloge n'est pas valide}, \'evitant un
d\'eclenchement \`a minuit par d\'efaut. \\

\textbf{Fonctionnement sans module d'horloge (DS3231).} & Mise en place d'une \textbf{horloge
logicielle de secours}, re-synchronis\'ee par le t\'el\'ephone \`a chaque connexion. \\

\textbf{Microcontr\^oleur incapable de notifier le t\'el\'ephone} directement. &
Adoption du \textbf{sondage par num\'ero d'\'ev\'enement}~: l'app interroge le distributeur
toutes les 3 s, sans technologie complexe (pas de WebSocket). \\

\textbf{\'Evolution du cahier des charges}~: connectivit\'e d'abord interdite, puis exig\'ee. &
Migration d'une approche \textbf{Bluetooth} vers une approche \textbf{WiFi/HTTP}, plus simple
\`a tester et compatible avec l'environnement Expo Go (sans module natif). \\
\bottomrule
\end{tabularx}

\subsection{Limites connues et pistes d'am\'elioration}

\begin{warnbox}
\begin{itemize}
  \item \textbf{Communication non chiffr\'ee et sans authentification}~: sur le r\'eseau
        local, toute personne connect\'ee pourrait commander le moteur. Acceptable pour un
        \textbf{prototype}, \`a s\'ecuriser pour une version finale.
  \item \textbf{Identifiants WiFi inscrits en clair} dans le code du firmware~: \`a
        externaliser dans un fichier de configuration prot\'eg\'e.
  \item \textbf{M\'emoire d'\'ev\'enements limit\'ee}~: un seul \'ev\'enement est conserv\'e
        entre deux interrogations. Si deux \'ev\'enements surviennent en moins de 3 secondes,
        le premier peut \^etre perdu.
\end{itemize}
\end{warnbox}

% ═════════════════════════════════════════════════════════════
\section*{Conclusion du chapitre}
\addcontentsline{toc}{section}{Conclusion du chapitre}

Ce chapitre a montr\'e que DisrtuCare repose sur une architecture \textbf{d\'ecoupl\'ee mais
coordonn\'ee}~: un distributeur \textbf{autonome} qui assure la fonction vitale (d\'elivrer la
bonne dose \`a la bonne heure) et une application \textbf{compagnon} qui enrichit
l'exp\'erience (rappels, suivi d'observance, t\'el\'ecommande). Leur int\'egration, via une
communication WiFi simple et robuste, permet qu'un \textbf{geste physique} du patient se
traduise automatiquement en \textbf{donn\'ee de suivi}. Les difficult\'es rencontr\'ees,
essentiellement mat\'erielles, ont toutes trouv\'e une solution, et les limites restantes
sont clairement identifi\'ees en vue d'une version aboutie.

\end{document}
